%% frontmatter/preface.tex —— 前言（主编维护）
\chapter*{第一版前言}
\addcontentsline{toc}{chapter}{第一版前言}
\markboth{第一版前言}{第一版前言}

\epigraph{我一直想为奶蛙作传。}{\itshape —— 第一版编者}

本书脱胎于一份在高校社团之间流传的小册子。

2026 年 9 月，一份名为《捧腹大笑思想与奶龙学概论》的二十页讲义出现在若干高校
奶龙同好会的群文件里。它用九节的篇幅记录了奶龙与奶蛙的来龙去脉，
用一种半认真半玩笑的口气，试图解释“为什么当下的大学生如此喜欢一只捧腹大笑的青蛙”。
它的写法是随笔式的：材料丰富、感情真挚、体例自由，末尾还留了一句
“鄙人奶疏学浅”。

那份小册子提出了一些真问题。它追问：为什么越是竞争激烈的顶尖高校，
越是流行这种无厘头的东西？它给出五种解释——情绪代偿、群体稀释、反精致、
学校社交、逆反认证。它甚至自觉地在末尾打趣说“扯不下去了”。
这些解释有洞见，但它们是\textbf{平行罗列}的：五条之间是什么关系？
哪一条可以被证伪？哪一条只是把现象换了个说法？——小册子没有回答，
或者说，它本来也没有打算回答。

本书打算回答。

\section*{本书对第一版做了三件事}

\textbf{第一，把“章节”改成“章—节—目”。}
第一版的九节在体例上是并列的，本书把它们重组为十二章，
每章内部遵循“提出问题—界定概念—展开分析—给出结论—留下习题”的教材结构。
历史材料归历史章，理论模型归理论章，实践案例归实践章，
使读者能看清“什么算是证据”“什么是解释”。

\textbf{第二，把“事实”与“说法”分开。}
第一版关于奶蛙起源的叙述，与今天社群考据的结果已有出入；
书中的部分年份（例如第二季剧集的上线时间）缺乏独立来源印证；
而关于奶蛙笑声音源、关于起源的那场考据过程，
在第一版成书时还远未结束。本书对每一条材料标注了可信度等级
（A 一手权威 / B 机构公开信息 / C 社群考据 / M 本书自建模型），
凡属推测一律写明“这是推测”。这不是对第一版的苛责——
恰恰相反，正是因为第一版把材料诚实地摊在了读者面前，
后来者才有可能把它们重新整理。

\textbf{第三，把“感想”变成“可检验的假设”。}
本书保留了第一版五重机制的直觉，但把它们重写为
“压力源—补偿行为—可观察指标—竞争性解释”的四列结构，
并明确指出：在没有问卷与访谈数据的情况下，
它们只是\textbf{待检验的假设}，不是结论（见第七章、第十一章）。

\section*{阅读本书需要什么}

不需要。本书假定读者没有任何传播学或社会学的预备知识，
所有术语都在首次出现处给出定义，所有外文概念都附原文。
如果你只想看乐子，第一章和第四章会满足你；
如果你想认真做一次文化研究，第十章和第十二章是必经之路。

\section*{关于本书的态度}

本书的研究对象是一种以“绷不住”为核心体验的文化现象。
写作这样的书有一个天然的危险：要么端着架子把人逗笑的东西分析得索然无味，
要么干脆放弃分析，跟着一起笑过去。本书的立场是——
\textbf{对现象保持幽默，对方法保持严肃}。
所以你会看到：章名可以开玩笑，图表必须是严谨的；
脚注可以吐槽，结论必须有出处。

至于这本书值不值得写、奶蛙值不值得研究——
我们把它交给读者的判断，以及时间。

\vspace{2em}
\begin{flushright}
主编\quad 忘川秋裤\\
\BookDate\quad 于清华园
\end{flushright}

\clearpage

\chapter*{编写说明与使用建议}
\addcontentsline{toc}{chapter}{编写说明与使用建议}
\markboth{编写说明与使用建议}{编写说明与使用建议}

\section*{一、本书的体例标识}

\begin{nlkeypoints}
  \begin{enumerate}[label=(\arabic*),leftmargin=2.4em]
    \item 每章开头以“本章要点”列出该章要回答的三个问题；
    \item 正文中\textbf{黑体术语}表示关键概念，完整释义见附录 B（术语表），
          并可通过书末“术语与人名索引”检索；
    \item 灰底青色框为\textbf{术语释义}，赭色框为\textbf{史料与考据}，
          紫色框为\textbf{研究方法提示}，红色框为\textbf{批判性反思}，
          绿色框为\textbf{案例}；
    \item 标有“假设 X.Y”（如“假设 7.5”）的编号条目，是本书提出但尚未经经验检验的命题，
          请勿当作已证结论引用；
    \item 每章末尾附“本章小结”“本章习题”“关键术语”。
  \end{enumerate}
\end{nlkeypoints}

\section*{二、可信度标注}

本书在涉及具体事实处使用如下简写：\conf{A} 表示可追溯到一手权威来源；
\conf{B} 表示机构或平台的公开信息；\conf{C} 表示社群考据或二手评述，
\textbf{只能作为线索}；\conf{M} 表示本书为教学目的构造的模型或示例数据。
所有 \conf{C} 级材料在正文中都以“据……”“有说法认为”的方式引出。

\section*{三、教学使用建议}

\begin{nlmethod}[教学使用建议]
  \textbf{16 学时方案}（通识选修）：第一、二、五、六、七、十章各 2 学时，
  第十一章 2 学时，其余章节作为阅读材料。\\
  \textbf{32 学时方案}（社团培训）：全书讲授，
  每章配一次小组作业，第十二章的方法与伦理训练单独安排 2 学时实践。\\
  \textbf{自主阅读}：推荐路径为第一章 → 第三章 → 第四章 → 第六章 → 第七章 →
  第十章 → 第十二章。
\end{nlmethod}
